\chapter{Einführung in die Medizin}
\label{chp:EIN}

\emph{Maintainer: Sebastian Gabriel}

\minitoc

\gDefinition{Die Medizin ist die Lehre von Gesundheit und
Krankheit sowie die Kunde der Gesunderhaltung, Heilung und
Linderung.}

% \paragraph*{Einteilung} Notfälle kann man mehr oder weniger willkürlich
% einteilen
% 
% \begin{itemize}
% 	\item Störungen des Bewu{\ss}tseins --  neurologische Notfälle
% 	\item pulmonale Notfälle
% 	\item Kreislaufstörungen
% 	\item Gefä{\ss}verschlüsse
% 	\item Abdominelle Notfälle
% 	\item Urologische Erkrankungen
% \end{itemize}

\section{Begriffe}

Die Kenntnis von Fachbegriffen in der Medizin und
Sanitätshilfe ist nicht deshalb wichtig, weil die Fachwörter
so gut klingen. Vielmehr dient die Fachsprache dazu, dass
verschiedene Personen miteinander kommunizieren können und
dabei sichergestellt ist, dass sie sich gegenseitig
verstehen.

\subsection{Klinik -- Diagnose -- Therapie}

\paragraph{Symptome} \index{Symptom} Krankheitserscheinungen
\paragraph{Komplikationen} \index{Komplikation}
möglicherweise auftretende Schwierigkeiten
\paragraph{Therapie} \index{Therapie} Ma{\ss}nahmen zur
Heilung oder Linderung von Krankheiten und Symptomen
\paragraph{Diagnose} \index{Diagnose}

Der Begriff Diagnose wird dauernd verwendet, dennoch ist er
nicht einfach zu erklären. Eine Diagnose gibt einem Zustand
eine Bezeichnung, das heißt einem Beschwerdebild wird ein
definiertes Krankheitsbild zugeordnet.

\begin{tabularx}{\linewidth}{XcX}\toprule
Beschwerden und Symptome des Patienten
&
{\color{darkBlue}\sffamily\bfseries $\rightarrow$~DIAGNOSE~$\leftarrow$}
&
Definiertes Krankheitsbild
\\\addlinespace
Patient
&
{\color{darkBlue}\sffamily\bfseries $\rightarrow$~DIAGNOSE~$\leftarrow$}
&
Lehrbuch
\\\addlinespace
Realität
&
{\color{darkBlue}\sffamily\bfseries $\rightarrow$~DIAGNOSE~$\leftarrow$}
&
Lehre
\\\bottomrule
\end{tabularx}

\gFazit{Die Diagnose ordnet die Beschwerden des Patienten einem definierten
Krankheitsbild zu.}

Mitunter kann diese Zuordnung einfach sein, sehr oft ist die Diagnosestellung
aber unsicher. Um die Unsicherheit dieser Diagnosen
anzuzeigen haben sich verschiedene Begriffe eingebürgert: 

\begin{description}
    \item[Verdachtsdiagnose] \index{Verdachtsdiagnose} Die
    Diagnose ist unsicher, man hat zwar den Verdacht in Richtung eines bestimmten Krankheitsbildes, andere ähnliche
    Krankheitsbilder können aber nicht mit hinreichender Sicherheit
    ausgeschlossen werden. Die endgültige Klärung überlässt man jemand
    anderen\footnote{z.B. dem Krankenhaus}. \par 
    Sanitäter erstellen i.d.R. nur Verdachtsdiagnosen.
     
    \item[Arbeitsdiagnose] \index{Arbeitsdiagnose}
    Arbeitsdiagnosen sind Verdachtsdiagnosen bei denen eine endgültige Klärung nicht möglich ist. Für die Behandlung wird die
    wahrscheinlichste Diagnose ausgewählt um \emph{``am Patienten arbeiten zu
    können''.}
    \item[Notfalldiagnose] \index{Notfalldiagnose} Eine
    Notfalldiagnose ist im Prinzip eine Arbeitsdiagnose. Der Patient ist dabei in einem Zustand bei dem
    \gEs{sofortiges Handeln} notwendig ist und die eigentliche Diagnose, die zu
    dieser Situation geführt hat, in dem Moment unwichtig ist. \par
    Ein Beispiel hierfür sind die Bewusstlosigkeit und der
    Kreislaufstillstand. Die Ursachen dafür sind erstmal
    unwichtig, mit der entsprechenden Behandlung der
    Notfalldiagnose muss sofort begonnen werden.
    \item['Status post' -- 'Zustand nach']
    \index{Status post}\index{Zustand
    nach}\index{St.p.|see{Status post}}
    \index{Z.n.|see{Zustand nach}} Abk. \gE{'St.p.'}, \gE{'Z.n.'}. Bezeichnet eine ``alte Diagnose'', d.h. eine Erkrankung oder eine Behandlung die einmal durchgemacht wurde. Z.B.:
    \begin{itemize}
        \item \gTt{Z.n. Blinddarm-Operation}: Der Patient
        hatte einmal eine Blinddarm-Operation.
        \item \gTt{Z.n. Herzinfarkt 2/2006}: Der Patient
        hatte im Februar 2006 einen Herzinfarkt.
    \end{itemize}
    \emph{``Umgangssprachlich''} wird oft auch der Umstand der zu einer
    Diagnose geführt hat mit \emph{``Z.n.''} angegeben. 
    \begin{itemize}
        \item \gTt{V.a. Schenkelhalsfraktur re., Z.n. Sturz}: Der
        Patient ist gestürzt und hat deshalb möglicherweise eine
        Schenkelhalsfraktur rechts.
    \end{itemize}
\end{description}
\paragraph{Differentialdiagnosen (DD)} \index{} Diagnosen,
welche bei den jeweiligen Symptomen auch
möglich, aber nicht so wahrscheinlich sind wie jene
Diagnose für welche man sich ``entscheidet''.

Beispiel: Patient mit Kopfschmerzen nach ausgiebigen
Alkoholgenuss am Vortag.
\begin{description}
    \item[Verdachtsdiagnose:] \gE{``Kater''}
    \item[Differentialdiagnosen:]
    Migraine, Hirnblutung, Schlaganfall,
    Hypoglykämie, Intoxikation, \ldots
\end{description}

% TODO ORTHO: Migraine? Migräne?

\paragraph{Diagnostik}\index{Diagnostik}

Als Diagnostik versteht man alle Bemühungen um eine Diagnose zu stellen, egal
ob man sich der eignen Sinne oder komplizierter Geräte
bedient.

\paragraph{Befund}\index{Befund}

Jede diagnostische Maßnahme hat (hoffentlich) ein Ergebnis, welches man als
\gT{Befund} bezeichnet. 

%\subsection{Medizinische Fachbereiche}

%\subsection{Andere Medizinische Berufe --- Ein kleines
%\emph{"`who is who"'}}

\subsection{Allgemeine Begriffe}

\begin{multicols}{2}


\begin{compactdesc}
    \item[akut] plötzlich; rascher Handlungsbedarf.
    \item[chronisch] dauerhaft, wiederkehrend.

\end{compactdesc}






\paragraph{Vorsilben}

\begin{compactdesc}
  \item[Hyper-] hoch, oben.
  \item[Hypo-] tief, unten.
  \item[Tachy-] schnell.
  \item[Brady-] langsam.
\end{compactdesc}


\paragraph{Stoffe}
\begin{compactdesc}
  \item[Amminosäuren] Bausteine von Proteinen.
  \item[Enzym] Hilfsstoff, welcher eine chemische Reaktion (Umwandlung von
  Stoffen in andere Stoffe) begünstigt.
  \item[Protein] Eiweißstoff. Proteine werden aus Amminosäuren
  zusammengesetzt.
  
\end{compactdesc}

\paragraph{Sonstige Begriffe}

\begin{compactdesc}
  \item[Bolus] Bisse
  \item[Cave] "`Achtung!"', "`Hüte Dich!"'
  \item[Obstruktion] Verlegung, Verengung. z.B. im Rahmen
  der \gAbk{COPD }(\gSee{topic:copd}).
  \item[Pathologie, pathologisch, Patho-] Lehre von den
  Krankheiten. Als pathologisch wird alles krankhafte
  bezeichnet. Die Vorsilbe "`Patho"' weißt auf etwas
  abnormales, krankhaftes hin.
  \item[Zyanose] Blaufärbung der Haut aufgrund eines mangelnden
  Sauerstoffgehalts des Blutes. Betroffen sind inbesonders die Finger, Lippen
  und das Gesicht.
\end{compactdesc}


\end{multicols}









\section{Wichtige Pathomechanismen}

\subsection{Entzündung}
\index{Entz\"undung}
\label{topic:entzuendung}

\gDefinition{Reaktion des K\"orpers welche charakterisiert
ist durch:}

\begin{enumerate}
    \item \gEs{R\"otung}
    \item \gEs{Schwellung}
    \item \gEs{\"Uberw\"armung}
    \item \gEs{Schmerz}
    \item \gEs{Funktionsverlust}
\end{enumerate}

\paragraph{Ausl\"oser}

\begin{itemize}
    \item Fremdk\"orper
    \item Krankheitserreger
    \item Mechanische Beanspruchung
    \item ...
\end{itemize}

\paragraph{Anh\"angsel
{\textquotedblleft{}\gT{-itis}{\textquotedblright}}}\index{itis@-itis}
Gastritis ${\rightarrow}$
{\quotedblbase}Magenentz\"undung{\textquotedblleft}, Hepatitis
${\rightarrow}$ Leberentz\"undung, Pankreatitis ${\rightarrow}$
Entz\"undung der Bauchspeicheldr\"use, Meningitis ${\rightarrow}$
Hirnhautentz\"undung, Rhinitis ${\rightarrow}$
{\quotedblbase}Schnupfen{\textquotedblleft}, Otitis ${\rightarrow}$
Ohrenentz\"undung, ...


\subsection{Ischämie}

% TODO Inhalt: Ischämie


\subsection{Tumore -- Neubildungen -- Krebs}

\gTakeNotes

Die meisten Zellen des menschlichen Körper teilen sich
regelmäßig. Diese Form der Zellvermehrung ist für das
Überleben des Menschen notwendig. Eine Kontrolle dieses
Vorgangs ist jedoch erforderlich um zu gewährleisten dass
die Körperzellen nicht wild zu wuchern beginnen. Der Körper
verfügt über verschiedene Mechanismen um die Zellteilung im
richtigen Maß zu halten. 

Wenn diese Mechanismen versagen kommt es zu einer
\gE{Entartung} und \gE{Wucherung}, meist in Form einer
\gEs{Raumforderung} und Schwellung (Tumor\footnote{Der Begriff \gE{"`Tumor"'} bedeutet ganz
allgemein "`Schwellung"'. Umgangssprachlich wird er oft
(fälschlicherweise) mit "`Krebs"' gleichgesetzt.}). Man
spricht dann von einer Neubildung \gZ{(Neoplasie)}, in
Arztbriefen oft mit "`{\ttfamily N.}"' abgekürzt. Der
medizinische Fachbereich der sich mit
Neubildungen (Krebserkrankungen) befasst ist die
\gE{Onkologie}.

Man unterscheidet zwischen gutartigen \gZ{(benignen)} und
bösartigen \gZ{(malignen)} Neubildungen. Gutartige
Neubildungen "`bleiben wo sie sind"' und wachsen vor sich
hin. Bösartige Neubildungen bilden Absiedelungen,
sog. \gE{Metastasen}, welche z.B. über die Blutbahn an
andere Stellen des Körper gelangen und dort zu wuchern
beginnen. 

Grundsätzlich können alle teilungsfähigen Zellen entarten.
Dies kann spontan geschehen oder durch andere Faktoren
(radioaktive Strahlung, chemische Substanzen) begünstigt
bzw. ausgelöst werden.

\gZ{Je nach Zellart unterscheidet man bei bösartigen
Wucherungen zwischen Karzinomen, Sarkomen und Blastomen.}

\begin{table}[H]
\gCaptionof{table}{Beispiele: Häufige Neubildungen.}

\begin{tabularx}{\linewidth}{XXl}
\gEs{Schlagwort}
&
\gEs{Beispiel}
&
\gTabHeading{\gZ{Im Arztbrief}}
\\\midrule
Lungenkrebs
&
Bronchialkarzinom
&
{\ttfamily\gZ{N. bronchii}}
\\\addlinespace
Brustkrebs
&
Mammakarzinom
&
{\ttfamily\gZ{N. mammae}}
\\\addlinespace
Darmkrebs
&
Rektumkarzinom
&
{\ttfamily\gZ{N. recti}}
\\\addlinespace
Blutkrebs
&
Leukämie, versch. Arten
&
{\ttfamily\gZ{ALL, AML, CML, \ldots}}
\\\bottomrule
\end{tabularx}
\end{table}

\paragraph{Was ist das Problem?}

Zusammengefasst ergeben sich drei wesentliche Probleme:

\begin{enumerate}
    \item \gEs{Neubildungen brauchen Raum.} Dieser Raum ist
    nicht endlos vorhanden. Je nach Lage können wichtige
    Strukturen beengt oder behindert werden (Hirn,
    Luftröhre, \ldots).
    \item \gEs{Funktionsverlust des gesunden Gewebes.}
    Durch den Befall mit (nutzlosen) Metastasen wird das
    gesunde Gewebe soweit durchwachsen und geschwächt dass
    es seine Funktion nicht mehr richtig ausführen kann
    (z.B. Knochenbefall).
    \item \gEs{Schmerz.} Der Befall verursacht oft starke
    Schmerzen. Krebspatienten benötigen oft eine
    dauerhafte, starke Schmerztherapie.
\end{enumerate}

\paragraph{Therapie}

Der Therapieerfolg ist sehr stark von der genauen Art der
Neubildung und vom jeweiligen Patienten abhängig und reicht
von sehr gut bis nicht vorhanden. 

Im wesentlichen kommen bei der Krebstherapie folgende
Methoden zum Einsatz:

\begin{description}
    \item[Chemotherapie] Einsatz von Medikamenten welche
    die Zellteilung unterdrücken bzw. Zellen absterben
    lassen.
    \item[Bestrahlung] zur Abtötung oder Vorbeugung von
    Geschwülsten.
    \item[Operation] Chirurgische Entfernung der
    Neubildung(en).
\end{description}

Diese Therapien sind für den Körper sehr anstrengend und
haben viele Nebenwirkungen, wie z.B. Haarausfall, Übelkeit,
Erbrechen, \gEs{Schwächung des Immunsystems}, Nervenschäden
u.v.a.m.

% \subsection{Ischämie}

% TODO INHALT-NEU: Ischämie 

% \section{Reserviert}






\section{Dokumentation -- Von der medizinischen Seite
betrachtet}
\index{Dokumentation!medizinische Aspekte}
\label{topic:dokumentation-medizinisch}

\emph{Vgl. auch \gSee{topic:dokumentation-rechtlich}}

\gParallel{
Die medizinische Dokumentation muss eine Geschichte erzählen. Eine
Geschichte die ein Dritter, unbeteiligter, verstehen
\gE{muss}. Dieser unbeteiligte Dritte soll allein anhand der
Schilderung in Ihrer Dokumentation fähig sein eine
Verdachtsdiagnose zu stellen. Weiters sollte er
nachvollziehen können, was mit dem Patient passiert ist.
Auch die Ereignisse davor oder die Einsatzsituation soll
für denjenigen verständlich sein (Unfallhergang,
Wohnsituation bei psychiatrischen Patienten, \ldots)
}{
\begin{center}
\includegraphics[width=0.5\linewidth]{./images/EIN/affe-sama-ambas-schein.jpg}
\end{center}
}



\minisec{Schlecht:}
\begin{quote}

{\sffamily Diagnose:} {\ttfamily\color{darkBlue} RQW}

{\sffamily Befunde:} {\sffamily RR}
{\ttfamily\color{darkBlue} 130/70,} {\sffamily HF}
{\ttfamily\color{darkBlue} 90} {\sffamily SpO2}
{\ttfamily\color{darkBlue} 98\% }

{\sffamily Anmerkungen: }{\ttfamily\color{darkBlue} -- --
--}
\end{quote}

\minisec{Besser:}

\begin{quote}
{\sffamily Diagnose:} {\ttfamily\color{darkBlue} RQW Stirn}

{\sffamily Befunde:} {\sffamily RR}
{\ttfamily\color{darkBlue} 130/70,} {\sffamily HF}
{\ttfamily\color{darkBlue} 90} {\sffamily SpO2}
{\ttfamily\color{darkBlue} 98\% }

{\sffamily Anmerkungen:} {\ttfamily\color{darkBlue} Pat.
fuhr mit Fahrrad und prallte mit Stirn gegen Ast., kein Sturz. Lt. Begleiter
keine B.losigkeit \underline{Traumacheck:} RQW Stirn,
ca 50\,ml Blutverlust sonst unauffällig. \underline{Neuro:}
grob neurologisch unauffällig, Bew. klar, Pat zeitl, örtl, zur Situation \& Person orientiert, Pupillen
mittelweit, Reakt. prompt und seitengleich, Kraft\&Gefühl
seitengl. in OE und UE. \underline{Anamn.:} \underline{Sympt:} Kopfschmerz seit
Aufprall. \underline{Krankh.:} keine bek. \underline{All.}
keine bek. \underline{Med.:} Aspirin heute früh. letzte
Tetanus-Impf. nicht bek. \underline{Ereignisse} s.o.
\underline{Letzte} Mahlzeit: heute Früh 2 Semmeln\gA{ KOCH:
? muss sein? GABS: ja, 1) wegen s-kamel, 2) wegen nüchtern } 

\underline{Maßnahmen:} Steriler Verband, nüchtern belassen,
Transport liegend mit erh OK.}
\end{quote}

\minisec{Psychiatrisch:}

\begin{quote}
{\sffamily Diagnose:} {\ttfamily\color{darkBlue}
Psychiatrische Erkrankung, V.a. Wahnvorstellungen, Unterbringung nach UbG}

{\sffamily Befunde:} {\sffamily RR}
{\ttfamily\color{darkBlue} 130/70,} {\sffamily HF}
{\ttfamily\color{darkBlue} 90} {\sffamily SpO2}
{\ttfamily\color{darkBlue} 98\% }

{\sffamily Anmerkungen:} {\ttfamily\color{darkBlue} Pat.
sitzend in Wohnung angetr., von Pol. mit Handschellen gefesselt. Pat. warf lt. Polizei Gegstände aus d Fenster
auf Straße, verletzte Passanten ($\rightarrow$ FLO-1). Pat.
gibt an sich geg. Geheimdienst gewehrt zu haben,
jetzt agitiert und unkooperativ. 

\underline{Traumacheck:} keine äußeren Anz f Verletzungen,
sonst nicht erhebbar (unkoop.). \underline{Neuro:} Pat.
ist wach, sonst nicht erhebbar

\underline{Anamn.:} \underline{Sympt:} keine Antwort
\underline{Krankh.:} lt. Spitalbrief Otto-Wagner-Sp. v.
1.4.08 paranoide Schizophrenie \underline{All.} k.A.
\underline{Med.:} nicht eruierbar, auf Tisch mehrere
angebr. Packungen Rohypnol \underline{Ereignisse} s.o.
\underline{Letzter Spitalsaufenthalt} k.A.

\underline{Wohnsituation:} heruntergekommen, vermüllt, Pat.
wohnt mit großem Hund, ganze Wohnung verkotet, bestialisch
stinkend. 

\underline{Maßnahmen:} Einweisung nach UbG, Transport in
Handschellen in Pol-begl. (Viktor 6, DNr. 4711). Protokoll
Amtsarzt ad. Spital.}
\end{quote}

% \section{Reserviert}
% \gA{Medizinische Fachbereiche}
% 
% \section{Reserviert}
% \gA{Medizinische Berufsgruppen}
% 
% \section{Reserviert}

\nocite{UlfigHistologie1}